\title{DeepSeekMath: Pushing the Limits of Mathematical Reasoning in Open Language Models}

\begin{abstract}
The structured intricacy inherent in mathematical reasoning makes it particularly difficult for language models to master. Here, we present DeepSeekMath~7B, a model obtained by extending the pre-training of DeepSeek-Coder-Base-v1.5 7B with 120B math-centric tokens, drawn from Common Crawl along with additional natural language and programming code data. Remarkably, DeepSeekMath~7B attains a score of 51.7\% on the advanced MATH benchmark without utilizing supplementary toolkits or ensemble methods, reaching the level of models like Gemini-Ultra and GPT-4. With a self-consistency evaluation over 64 generations, DeepSeekMath~7B achieves 60.9\% accuracy on MATH. The heightened mathematical reasoning of DeepSeekMath~is a consequence of two primary contributions: First, an extensive exploitation of open web resources through a carefully crafted data selection framework; second, the introduction of Group Relative Policy Optimization (GRPO), an adaptation of Proximal Policy Optimization (PPO), which simultaneously improves mathematical reasoning and reduces PPO's memory consumption.
\end{abstract}

\section{Introduction}

The emergence of large language models (LLMs) has fundamentally transformed how artificial intelligence addresses mathematical reasoning tasks, leading to notable progress on quantitative reasoning benchmarks \citep{MATH} and geometry-focused challenges \citep{trinh2024solving}. Additionally, such models have demonstrated valuable utility in assisting humans to tackle sophisticated mathematical questions \citep{tao}. Nonetheless, leading systems like GPT-4 \citep{gpt4} and Gemini-Ultra \citep{gemini} remain inaccessible to the public, resulting in a substantial gap in performance between these proprietary models and those currently available as open source.

This work presents DeepSeekMath, a specialized language model designed for mathematics, which demonstrates superior mathematical reasoning compared to other open-source models and nearly matches GPT-4’s performance on academic benchmarks. Central to our approach is the development of the DeepSeekMath Corpus, a comprehensive, high-quality pre-training dataset containing 120B mathematical tokens. To assemble this corpus, we leverage the Common Crawl (CC) and apply a fastText-based classifier \citep{joulin2016fasttext}. In the initial stage, the classifier is built with positive samples sourced from OpenWebMath \citep{openwebmath} and a broad array of web content as negative samples. After this phase, the classifier is used to identify additional mathematical content from CC, with subsequent refinement performed through human annotation. The improved dataset is then used to retrain the classifier, further boosting its extraction capabilities. Empirical assessments reveal that this large-scale corpus provides robust quality, as evidenced by DeepSeekMath-Base 7B reaching 64.2\% on GSM8K \citep{gsm8k} and 36.2\% on the advanced MATH benchmark \citep{MATH}, surpassing the performance of Minerva 540B \citep{minerva}. Notably, because the DeepSeekMath Corpus includes multilingual data, we also observe enhanced outcomes on Chinese mathematical assessments \citep{wei2023cmath,agieval}. We suggest that our methodologies in curating and processing mathematical data offer a valuable starting foundation for further advancements, with substantial opportunities for future improvement.

DeepSeekMath-Base is initialized using DeepSeek-Coder-Base-v1.5 7B \citep{deepseek-coder}, as empirical evidence suggests that leveraging a model pre-trained on code yields superior outcomes relative to initializing from a generic large language model (LLM). In addition, mathematical training contributes not only to mathematical proficiency but also improves model performance on the MMLU \citep{mmlu} and BBH benchmarks \citep{bbh}, suggesting that this process enhances both quantitative and general reasoning skills.

Subsequent to pre-training, DeepSeekMath-Base undergoes mathematical instruction tuning, incorporating datasets designed for chain-of-thought \citep{cot}, program-of-thought \citep{pot,pal}, and reasoning augmented with external tools \citep{tora}. The resultant model, DeepSeekMath-Instruct 7B, surpasses all other instruction-tuned 7B models and demonstrates competitiveness with the open-source 70B counterparts.

We further propose Group Relative Policy Optimization (GRPO), a modified reinforcement learning (RL) algorithm based on Proximal Policy Optimization (PPO) \citep{schulman2017proximal}. GRPO eliminates the need for a critic network by estimating the baseline using group-based scores, thereby markedly decreasing computational requirements during training. Utilizing only a subset of English instruction-tuning data, GRPO achieves notable performance gains over the strong DeepSeekMath-Instruct model across both in-domain (GSM8K: 82.9\% $\rightarrow$ 88.2\%, MATH: 46.8\% $\rightarrow$ 51.7\%) and out-of-domain tasks (e.g., CMATH: 84.6\% $\rightarrow$ 88.8\%) in the RL phase. Additionally, we present a unified framework to interpret various methodologies such as Rejection Sampling Fine-Tuning (RFT) \citep{yuan2023scaling}, Direct Preference Optimization (DPO) \citep{dpo}, PPO, and GRPO. Within this unified context, all methods can be categorized as direct or simplified forms of RL. We perform thorough empirical studies—including comparisons between online and offline training, outcome versus process supervision, as well as single-turn and iterative RL—to analyze fundamental components of this unified framework. Finally, we elucidate the mechanisms through which our RL approach enhances instruction-tuned model performance and discuss future research directions to further optimize RL methods within this paradigm.

\subsection{Contributions}
This work introduces advancements in large-scale math pre-training and presents an investigation into reinforcement learning approaches.

\textbf{Math Pre-Training at Scale}

\begin{itemize}[topsep=0pt]
    \item Our study demonstrates that Common Crawl, despite its public availability, holds substantial mathematical information. Leveraging a carefully constructed filtering and selection process, we build the DeepSeekMath Corpus—a comprehensive dataset of 120B tokens derived from web content with a mathematical focus. This collection significantly exceeds previous efforts: it is nearly 7 times larger than the math web dataset utilized by Minerva \citep{minerva} and approximately 9 times the volume of the recent OpenWebMath \citep{openwebmath}.

    \item The pre-trained DeepSeekMath-Base 7B model we introduce attains performance on par with Minerva 540B \citep{minerva}, suggesting that model size alone does not dictate success in mathematical reasoning tasks. Instead, smaller architectures, when trained on carefully curated, high-quality data, can exhibit comparable or superior capabilities.

    \item We report key insights from our math-focused training investigations. Specifically, exposing the model to code training before mathematical pre-training leads to improved proficiency in solving math tasks, both with and without supplementary tools. These results contribute to the ongoing discussion regarding the impact of code pre-training on reasoning: our evidence supports its positive effect, at least within mathematical domains.

    \item Contrary to prevalent practice, supplementing training with arXiv math papers—despite their frequent inclusion in mathematics-oriented projects—yields no observable performance gain on the suite of mathematical benchmarks we employ.
\end{itemize}

\textbf{Exploration and Analysis of Reinforcement Learning}

\begin{itemize}[topsep=0pt]
    \item We present Group Relative Policy Optimization (GRPO), a novel reinforcement learning approach designed for both efficiency and efficacy. GRPO eliminates the need for a critic network by deriving the baseline directly from group-based score comparisons, thereby substantially lowering computational demands relative to Proximal Policy Optimization (PPO).
    \item Our results show that GRPO provides substantial improvements to the performance of the instruction-tuned DeepSeekMath-Instruct model, leveraging only instruction-tuning data. Additionally, we observe that GRPO leads to gains in generalization, enhancing out-of-domain outcomes during reinforcement learning.
    \item We propose a unified conceptual framework to contextualize various techniques, such as RFT, DPO, PPO, and GRPO. Our empirical study spans multiple dimensions—including comparisons of online and offline training strategies, analysis of outcome- versus process-based supervision, and exploration of both single-step and iterative reinforcement learning—to rigorously dissect the fundamental factors within this unified framework.
    \item Building on this integrative perspective, we analyze the underlying mechanisms that contribute to the success of reinforcement learning, and highlight several promising avenues for further advancing reinforcement learning methodologies for large language models (LLMs).
\end{itemize}

\subsection{Summary of Evaluations and Metrics}

\begin{itemize}[topsep=0pt]
    \item \textbf{English and Chinese Mathematical Reasoning}: We perform thorough evaluations of our models across a range of English and Chinese benchmarks, which encompass mathematical tasks from the elementary to the collegiate level. The English datasets include GSM8K \citep{gsm8k}, MATH \citep{MATH}, SAT \citep{llemma}, OCW Courses \citep{minerva}, and MMLU-STEM \citep{mmlu}; for Chinese, we utilize MGSM-zh \citep{mgsm}, CMATH \citep{wei2023cmath}, Gaokao-MathCloze \citep{agieval}, and Gaokao-MathQA \citep{agieval}. The evaluations assess both the models’ capabilities for producing complete, self-explanatory text-based solutions without auxiliary tools, as well as their proficiency in problem solving using Python.
    
    On the English benchmarks, DeepSeekMath-Base demonstrates competitive results in comparison to the proprietary Minerva 540B \citep{minerva}, and exceeds the performance of all other publicly available base models (such as Mistral 7B \citep{mistral} and Llemma-34B \citep{llemma}), whether or not they benefit from mathematical pre-training. The performance gap is often substantial. Importantly, DeepSeekMath-Base also achieves leading results on Chinese evaluation sets, a likely outcome of our methodology that, unlike prior efforts \citep{minerva, llemma}, incorporates high-quality non-English mathematical data in addition to English sources. Following instruction-based mathematical tuning and subsequent reinforcement learning, DeepSeekMath-Instruct and DeepSeekMath-RL achieve notable advancements, most prominently achieving over 50\% accuracy on the highly challenging MATH benchmark—marking a first for open-source models.
\end{itemize}

\item \textbf{Formal Mathematics}: We assess the capabilities of DeepSeekMath-Base on the informal-to-formal theorem proving benchmark introduced by \citep{dsp_proof}, using miniF2F \citep{minif2f} with Isabelle \citep{isabelle} serving as the designated proof assistant. DeepSeekMath-Base exhibits notable performance in few-shot autoformalization scenarios.
\item \textbf{Natural Language Understanding, Reasoning, and Code}: To thoroughly evaluate the models’ general proficiency in comprehension, reasoning, and programming, DeepSeekMath-Base is benchmarked on the Massive Multitask Language Understanding (MMLU) dataset \citep{mmlu}, featuring 57 diverse multiple-choice subject areas. We also test on BIG-Bench Hard (BBH) \citep{bbh}, which comprises 23 complex, primarily multi-step reasoning tasks, as well as HumanEval \citep{codex} and MBPP \citep{mbpp}, both established code evaluation benchmarks. Results demonstrate that mathematical pre-training enhances both language understanding and reasoning abilities.

\section{Math Pre-Training}

\subsection{Data Collection and Decontamination} 

This section details the methodology for assembling the DeepSeekMath Corpus from Common Crawl. As illustrated in Figure \ref{fig:data_collect}, we introduce an iterative data acquisition pipeline for methodically collecting a large-scale mathematical dataset from Common Crawl, initializing the process with a seed corpus such as a limited but high-quality set of mathematical data. This procedure can be generalized to other fields, including code-related datasets.

Initially, we select OpenWebMath \citep{openwebmath}—a curated resource of high-quality mathematical texts from the web—as the foundational seed corpus. Using this dataset, a fastText model \citep{joulin2016fasttext} is trained to retrieve web pages that resemble those in OpenWebMath. To do this, 500,000 randomly chosen entries from the seed serve as positive samples, and an additional 500,000 web pages from Common Crawl function as negatives. Model training utilizes an open-source fastText implementation\footnote{\url{https://fasttext.cc}}, with parameter settings as follows: a vector size of 256, learning rate of 0.1, a maximum n-gram length of 3, minimum word frequency set to 3, and 3 epochs. To streamline Common Crawl, URL-based deduplication along with near-duplicate removal techniques are applied, narrowing the dataset to 40B HTML pages. The fastText model then selects candidate mathematical pages from this deduplicated subset. Subsequently, the filtered content is ranked based on fastText scores, and only the highest-scoring documents are retained to ensure data quality. Pre-training experiments are carried out using datasets of varying sizes—top 40B, 80B, 120B, and 160B tokens—to assess retention thresholds. In the pipeline’s initial round, the top 40B tokens are chosen for further use.

Following the initial round of data collection, a significant number of mathematical web pages remain uncollected. This shortfall primarily stems from the limited diversity in the positive examples used to train the fastText model. To address this limitation, we incorporate additional mathematical web domains into the initial corpus to enhance the breadth of examples for model training. We first partition the entire Common Crawl dataset by unique domains, each defined as web pages sharing a base URL. For every domain, we compute the fraction of web pages already collected during the initial phase. Domains in which more than 10\% of their pages are retrieved are categorized as mathematics-related (e.g., \url{mathoverflow.net}). 

Next, we manually identify and label URLs within these candidate domains that point to mathematical resources (for instance, \url{mathoverflow.net/questions}). Web pages associated with these mathematical URLs but missing from the initial corpus are subsequently included to further augment the seed data. Through this expanded positive sample pool, the refined fastText model is able to recover more mathematical documents in each subsequent iteration. After four data collection cycles, the resulting set comprises 35.5M mathematical web pages with a total of 120B tokens. At the fourth round, we observe that nearly 98\% of the documents were already obtained during the third pass, prompting us to halt further data collection.

In order to prevent data leakage into evaluation benchmarks, we adopt the approach described by \cite{deepseek-coder} to eliminate any web pages containing questions or answers from prominent English mathematical benchmarks—such as GSM8K \citep{gsm8k} and MATH \citep{MATH}—as well as Chinese benchmarks, including CMATH \citep{wei2023cmath} and AGIEval \citep{agieval}. The filtering protocol entails discarding any text segment that matches a 10-gram substring from any of the benchmark datasets. For reference texts with fewer than 10 grams but no fewer than 3 grams, we apply a strict exact-match filter to identify and remove contaminated documents from our math training set.

\subsection{Validation of the DeepSeekMath~Corpus Quality}

To assess the quality of the DeepSeekMath~Corpus in comparison to newly available mathematics corpora, we conduct pre-training evaluations on several established datasets:

\begin{itemize}[topsep=0pt]
    \item \textbf{MathPile} \citep{mathpile}: a comprehensive corpus comprising 8.9B tokens aggregated from a variety of sources, including textbooks, Wikipedia, ProofWiki, CommonCrawl, StackExchange, and predominantly arXiv (contributing over 85\%);
    \item \textbf{OpenWebMath} \citep{openwebmath}: extracted from CommonCrawl, this dataset contains 13.6B tokens filtered to capture mathematical content;
    \item \textbf{Proof-Pile-2} \citep{llemma}: built from the combination of OpenWebMath, AlgebraicStack (which includes 10.3B tokens of mathematical code), and arXiv papers (28.0B tokens); following \cite{llemma}, experiments on Proof-Pile-2 adopt an arXiv:Web:Code data ratio of 2:4:1.
\end{itemize}

\subsubsection{Training Setting}
\label{sec:quality-policy}
We conduct specialized mathematical training using a general-purpose language model comprising 1.3B parameters, following the same architecture as the DeepSeek LLMs \citep{deepseek-llm}, henceforth referred to as DeepSeek-LLM 1.3B. Each mathematical corpus is used to train a separate model, with each undergoing training for 150B tokens. All training procedures are executed via the lightweight and efficient HAI-LLM framework \citep{haillm}. In line with DeepSeek LLM training strategies, we utilize the AdamW optimizer \citep{adamW} (setting $\beta_1=0.9$, $\beta_2=0.95$, and $\mathrm{weight\_decay}=0.1$) and implement a multi-phase learning rate schedule: the learning rate ascends to its maximum over the first 2,000 warmup steps, is reduced to 31.6\% of its peak by the 80\% completion mark, and then further lowered to 10.0\% of the maximum by 90\% progress. The learning rate's upper bound is 5.3e-4. We employ a batch size of 4M tokens, paired with a context length of 4K.

\subsubsection{Evaluation Results}

\textbf{The DeepSeekMath~Corpus demonstrates superior quality, comprehensive multilingual coverage in mathematics, and stands as the largest corpus of its kind.}

\begin{itemize}[topsep=0pt]
    \item \textbf{High-quality}:
    We measure downstream task performance across 8 mathematical evaluation benchmarks using few-shot chain-of-thought prompting \cite{cot}. Table \ref{tab:corpora_comparison} indicates that the model trained with the DeepSeekMath~Corpus outperforms others. Figure \ref{fig:corpora_comparisons} further highlights that DeepSeekMath-trained models achieve better results than those trained on Proof-Pile-2 after 50B tokens (equivalent to one epoch of Proof-Pile-2), suggesting a higher average sample quality within DeepSeekMath.
    \item \textbf{Multilingual}:
    The DeepSeekMath~Corpus includes material in various languages, with English and Chinese comprising the largest shares. According to Table \ref{tab:corpora_comparison}, pretraining with DeepSeekMath~Corpus improves mathematical reasoning in both English and Chinese, whereas previous corpora, which mostly focus on English, offer limited or negative transfer to Chinese mathematical reasoning.
    \item \textbf{Large-scale}:
    The scale of DeepSeekMath~Corpus surpasses that of other available mathematical corpora by several multiples. Figure \ref{fig:corpora_comparisons} shows that when DeepSeek-LLM 1.3B is trained on DeepSeekMath~Corpus, it exhibits more rapid and sustained performance gains. In contrast, the baseline datasets, being significantly smaller and repeatedly used throughout training, yield models whose performance quickly plateaus.
\end{itemize}

\subsection{Training and Evaluating DeepSeekMath-Base 7B}

This section details the development of DeepSeekMath-Base 7B, a foundational model exhibiting advanced reasoning proficiency, with particular strength in mathematical domains. The initial parameters of our model are sourced from DeepSeek-Coder-Base-v1.5 7B \citep{deepseek-coder}, after which it undergoes training over a corpus encompassing 500B tokens. The composition of the training dataset is as follows: 56\% DeepSeekMath Corpus, 4\% AlgebraicStack, 10\% arXiv content, 20\% Github code repositories, and 10\% natural language content extracted from Common Crawl in both English and Chinese. The training protocol generally adheres to the configuration outlined in Section \ref{sec:quality-policy}, with the exception that the peak learning rate is adjusted to 4.2e-4, and the batch size is set to 10 million tokens.

We systematically evaluate DeepSeekMath-Base 7B's mathematical reasoning capabilities by measuring its proficiency in generating complete, self-sufficient solutions, utilizing computational tools for problem-solving, and performing formal theorem verification. Additionally, we broaden our analysis to explore the model's performance in natural language understanding, complex reasoning, and programming to deliver a comprehensive overview of its base competencies.

\paragraph{Mathematical Problem Solving with Step-by-Step Reasoning}

To assess DeepSeekMath-Base's ability in stepwise mathematical problem solving, we employ few-shot chain-of-thought prompting \citep{cot} across a suite of eight benchmarks administered in both English and Chinese. These benchmarks include those focused on quantitative reasoning, such as GSM8K \citep{gsm8k}, MATH \citep{MATH}, and CMATH \citep{wei2023cmath}, as well as multiple-choice datasets like MMLU-STEM \citep{mmlu} and Gaokao-MathQA \citep{agieval}, spanning a wide array of mathematical disciplines from the elementary level up to college-level subject matter.

Table \ref{tab:base_math_cot} presents the results, showing that DeepSeekMath-Base 7B consistently achieves the highest scores across all eight evaluation benchmarks when compared to other open-source base models, such as the popular Mistral 7B \citep{mistral} and the recently launched Llemma 34B \citep{llemma}, which is trained on the Proof-Pile-2 dataset \citep{llemma}. Particularly noteworthy is its performance on the competition-level MATH dataset, where DeepSeekMath-Base exceeds all existing open-source base models by a margin greater than 10\% in absolute terms and even surpasses Minerva 540B \citep{minerva}—a closed-source base model that is 77 times larger, built upon PaLM \citep{palm}, and additionally fine-tuned with mathematical corpora.

\paragraph{Mathematical Problem Solving with Tool Use}

To assess capabilities in program-supported mathematical reasoning, we examine model performance on GSM8K and MATH using few-shot program-of-thought prompting strategies \citep{pot,pal}. For each test item, models generate a Python script leveraging available libraries such as \textit{math} and \textit{sympy} for handling complex calculations; the executed program's output is taken as the predicted solution. Results displayed in Table \ref{tab:base_math_pot_proof} indicate that DeepSeekMath-Base 7B establishes a new state-of-the-art among open-source models, outperforming the previously leading Llemma 34B.

\paragraph{Formal Mathematics}

Automating formal proof generation offers increased rigor, dependability, and productivity in mathematical verification, a topic that has recently garnered significant research interest. We test DeepSeekMath-Base 7B on the informal-to-formal proof task introduced by \citep{dsp_proof}, where the goal is to generate a formal proof given an informal statement, its formal equivalent, and an informal proof sketch. Our evaluation utilizes miniF2F \citep{minif2f}, a benchmark dedicated to formal mathematics at the Olympiad level; we apply few-shot prompting to direct the model to output formal Isabelle proofs. Consistent with \cite{dsp_proof}, we prompt the model to generate proof outlines and then employ the standard automated theorem prover Sledgehammer \citep{sledgehammer} to elaborate any omitted parts. Table \ref{tab:base_math_pot_proof} demonstrates that DeepSeekMath-Base 7B delivers robust performance in automatic formalization of proofs.

\paragraph{Natural Language Understanding, Reasoning, and Code}

We further examine model competence in natural language understanding (via MMLU \citep{mmlu}), reasoning (using BBH \citep{bbh}), and programming (on HumanEval \citep{codex} and MBPP \citep{mbpp}). According to Table \ref{tab:understanding_reasoning_code}, DeepSeekMath-Base 7B achieves notable gains in both MMLU and BBH relative to its predecessor, DeepSeek-Coder-Base-v1.5 \citep{deepseek-coder}, suggesting that its mathematical training yields broader improvements for language comprehension and logical reasoning tasks. Moreover, continual training on code-related data ensures that DeepSeekMath-Base 7B sustains the strong performance of DeepSeek-Coder-Base-v1.5 on the coding benchmarks. Across the board, DeepSeekMath-Base 7B substantially outperforms the widely used Mistral 7B \citep{mistral} on the trio of reasoning and programming evaluations.

\section{Supervised Fine-Tuning}

\subsection{SFT Data Curation}

We assemble a mathematical instruction-tuning corpus comprising problems in both English and Chinese, sourced from various mathematical disciplines and spanning a range of difficulty levels. Each problem is accompanied by a corresponding solution, formatted according to chain-of-thought (CoT) \citep{cot}, program-of-thought (PoT) \citep{pot,pal}, or tool-integrated reasoning frameworks \citep{tora}. The dataset contains a total of 776K training examples.

\begin{itemize}[topsep=0pt]
    \item \textbf{English mathematical datasets}: Our process includes the annotation of GSM8K and MATH problem sets with tool-integrated solution methods. Additionally, we incorporate a subset of MathInstruct \citep{MathInstruct} as well as the training partition of Lila-OOD \citep{lila}, where solutions follow either CoT or PoT styles. The English dataset encompasses numerous mathematical domains such as algebra, probability, number theory, calculus, and geometry.
    \item \textbf{Chinese mathematical datasets}: For Chinese-language data, we gather K-12 math problems distributed over 76 sub-topics—examples include linear equations—each provided with solutions represented in both CoT and tool-integrated reasoning styles.
\end{itemize}

\subsection{Training and Evaluating DeepSeekMath-Instruct 7B}

This section details the fine-tuning process for DeepSeekMath-Instruct 7B, an instruction-tuned variant derived from DeepSeekMath-Base targeting mathematical problem solving. We prepare the model input by concatenating training instances at random, up to a cap of 4K tokens per sequence. Training is conducted for 500 steps, using a batch size of 256, and employing a fixed learning rate of 5e-5.

To assess mathematical proficiency, we measure model performance—with and without access to tools—across four quantitative reasoning benchmarks in both English and Chinese. Comparative evaluation is carried out against contemporary leading models:

\begin{itemize}[topsep=0pt]
    \item \textbf{Closed-source models} under consideration comprise: (1) members of the GPT series, with GPT-4 \citep{gpt4} and GPT-4 Code Interpreter~\footnote{\url{https://openai.com/blog/chatgpt-plugins\#\#code-interpreter}} representing the current state of the art, (2) Gemini Ultra and Pro \citep{gemini}, (3) Inflection-2 \citep{inflection-2}, (4) Grok-1~\footnote{\url{https://x.ai/model-card}}, as well as recently launched models from Chinese enterprises including (5) Baichuan-3~\footnote{\url{https://www.baichuan-ai.com}}, and (6) the latest release in the GLM family, GLM-4~\footnote{\url{https://open.bigmodel.cn/dev/api\#glm-4}} \citep{glm}.
\end{itemize}

These models are designed for broad applicability and typically have been subject to several alignment processes. 
\item \textbf{Open-source models} consist of: general-purpose LLMs such as (1) DeepSeek-LLM-Chat 67B \citep{deepseek-llm}, (2) Qwen 72B \citep{qwen}, (3) SeaLLM-v2 7B \citep{seallm}, and (4) ChatGLM3 6B \citep{chatglm3}; in addition to variants tailored for mathematical tasks, including (5) InternLM2-Math 20B~\footnote{\url{https://github.com/InternLM/InternLM-Math}}, which extends InternLM2 with further mathematics-oriented training and instruction tuning, (6) Math-Shepherd-Mistral 7B, leveraging PPO training \citep{schulman2017proximal} with a process-supervised reward structure based on Mistral 7B \citep{mistral}, (7) the WizardMath series \citep{wizardmath}, which enhances mathematical reasoning for Mistral 7B and Llama-2 70B \citep{llama2} through evolve-instruct (a variant of instruction tuning utilizing AI-generated prompts) and PPO training, predominantly using problems from GSM8K and MATH, (8) MetaMath 70B \citep{metamath}, achieved by further fine-tuning Llama-2 70B with enriched versions of GSM8K and MATH, (9) ToRA 34B \cite{tora}, representing CodeLlama 34B further fine-tuned for tool-augmented mathematical reasoning, and (10) MAmmoTH 70B \citep{MathInstruct}, involving Llama-2 70B trained with MathInstruct for better alignment with mathematical instruction-following tasks.
\end{itemize}

According to Table \ref{tab:sft_rl_math}, in scenarios prohibiting the use of external tools, DeepSeekMath-Instruct 7B exhibits impressive capabilities in stepwise mathematical reasoning. On the MATH benchmark, which features competition-level problems, this model outperforms not only every open-source alternative but also the majority of commercial systems (such as Inflection-2 and Gemini Pro) by an absolute margin of at least 9\%. This superior performance holds even relative to models with substantially greater parameter counts (e.g., Qwen 72B) or those improved specifically via reinforcement learning on math-related tasks (e.g., WizardMath-v1.1 7B). While DeepSeekMath-Instruct demonstrates performance on par with top Chinese proprietary models like GLM-4 and Baichuan-3 on the MATH benchmark, it is still surpassed by the leading proprietary systems GPT-4 and Gemini Ultra.

When evaluation conditions permit the combination of natural language reasoning with programming-based tool use, DeepSeekMath-Instruct 7B achieves nearly 60\% accuracy on the MATH dataset—outperforming all other open-source models. Across additional benchmarks, this model remains competitive with DeepSeek-LLM-Chat 67B, the former state-of-the-art, despite being ten times smaller.

\section{Reinforcement Learning}

\subsection{Group Relative Policy Optimization} Reinforcement learning (RL) has consistently demonstrated utility in augmenting the mathematical reasoning competencies of LLMs subsequent to the Supervised Fine-Tuning (SFT) phase \citep{wang2023math,wizardmath}. Here, we present Group Relative Policy Optimization (GRPO), our RL method developed for both efficacy and efficiency.

\subsubsection{From PPO to GRPO} Proximal Policy Optimization (PPO) \citep{schulman2017proximal} serves as a widely-adopted actor-critic algorithm during the RL fine-tuning of LLMs \citep{ouyang2022training}. Its objective is to refine the LLM policy by maximizing the surrogate function defined as follows:

\begin{equation}
\footnotesize
    \mathcal{J}_{PPO}(\theta) = \mathbb{E}{[q \sim P(Q), o \sim \pi_{\theta_{old}}(O|q)]} \frac{1}{|o|} \sum_{t=1}^{|o|} \min \left[ \frac{\pi_\theta(o_{t} | q, o_{<t})}{\pi_{\theta_{old}}(o_{t} | q, o_{<t})} A_{t}, \text{clip} \left( \frac{\pi_\theta(o_{t} | q, o_{<t})}{\pi_{\theta_{old

Here, $\pi_{\theta}$ and $\pi_{\theta_{old}}$ represent the current and previous versions of the policy models, respectively, and $q$ and $o$ denote questions and outputs, sampled either from the question dataset or generated by the previous policy $\pi_{\theta_{old}}$. The hyper-parameter $\varepsilon$ is employed for clipping, a key element in PPO that helps maintain stable training dynamics. The term $A_t$ stands for the advantage, estimated via Generalized Advantage Estimation (GAE) \citep{gae}, which utilizes the sequence of rewards $\{r_{\ge t}\}$ in conjunction with a learned value function $V_{\psi}$. This necessitates training a value function in tandem with the policy model in PPO, and to counteract excessive optimization toward the reward model, it is common to introduce a per-token KL penalty, computed with respect to a reference model, at each step in the reward formulation \citep{ouyang2022training}:

\begin{equation}
    r_{t} = r_\varphi(q, o_{\le t}) - \beta \log\frac{\pi_{\theta}(o_{t}|q, o_{<t})}{\pi_{ref}(o_{t}|q, o_{<t})},
\label{eq:PPO-reward}
\end{equation}

where $r_\varphi$ is the output from the reward model, $\pi_{ref}$ refers to the reference model, typically the original SFT model, and $\beta$ acts as the scaling factor for the KL regularization.

Because the value function utilized in PPO is often similar in scale and complexity to the policy model itself, this approach can result in significant computational and memory requirements. Furthermore, throughout the reinforcement learning process, the value function serves as a baseline for the advantage calculation to minimize variance. However, within the context of large language models, it is generally the case that the reward model supplies feedback only for the final token, making the learning of a value function that is precise at each timestep more challenging. To overcome these issues, we introduce Group Relative Policy Optimization (GRPO), depicted in Figure \ref{fig:grpo}, which eliminates the need for an explicit value function, as in PPO. Instead, GRPO uses the average reward across multiple output samples generated for the same input question as a baseline. Specifically, for a given question $q$, GRPO samples a set of responses $\{o_1, o_2, \cdots, o_G\}$ using the old policy $\pi_{\theta_{old}}$ and then optimizes the policy by maximizing the following objective:

\begin{equation}
\footnotesize
\begin{split}
    \mathcal{J}_{GRPO}(\theta) &= \mathbb{E}{[q \sim P(Q), \{o_i\}_{i=1}^G \sim \pi_{\theta_{old}}(O|q)]}  \\
    & \frac{1}{G}\sum_{i=1}^G\frac{1}{|o_i|} \sum_{t=1}^{|o_i|} \left\{ \min \left[ \frac{\pi_\theta(o_{i,t} | q, o_{i,<t})}{\pi_{\theta_{old}}(o_{i,t} | q, o_{i,<t})} \hat{A}_{i,t}, \text{clip} \left( \frac{\pi_\theta(o_{i,t} | q, o_{i,<t})}{\pi_{\theta_{old}}(o_{i,t} | q, o_{i,<t})}, 1 - \varepsilon, 1 + \varepsilon \right)  \hat{A}_{i,t} \right] - \beta \mathbb{D}_{KL}\left[\pi_{\theta} || \pi_{ref}\right]\right\} ,
\end{split}
\label{eq:GRPO-obj}
\end{equation}

In this expression, $\varepsilon$ and $\beta$ remain as hyper-parameters, while $\hat{A}_{i,t}$ is the group-relative advantage, calculated exclusively from the rewards of the outputs within the same group—a method whose specifics are elaborated in the subsequent subsections. This group-based approach for calculating advantages in GRPO naturally corresponds to the comparative structure of reward model datasets, since these are typically built from paired comparisons of answers to the same prompt. Additionally, rather than incorporating the KL penalty directly into the reward, GRPO introduces a regularization term by adding the KL divergence between the current policy and the reference policy directly to the loss, thus streamlining the computation of $\hat{A}_{i,t}$.

\begin{algorithm}[t]
  \small
  \caption{Iterative Group Relative Policy Optimization}
  \textbf{Input} initial policy model $\pi_{\theta_{\text{init}}}$; reward models $r_\varphi$; task prompts $\mathcal{D}$; 
  hyperparameters $\varepsilon$, $\beta$, $\mu$
  \begin{algorithmic}[1]
    \State Set policy model $\pi_\theta \leftarrow \pi_{\theta_{\text{init}}}$
    \For{iteration = 1, \dots, I}
       \State Assign reference model $\pi_{ref} \leftarrow \pi_{\theta}$
      \For{step = 1, \dots, M}
      \State Draw a batch $\mathcal{D}_b$ from $\mathcal{D}$
      \State Copy current policy model $\pi_{\theta_{old}} \leftarrow \pi_{\theta}$ 
      \State For each question $q \in \mathcal{D}_b$, generate $G$ outputs $\{o_i\}_{i=1}^G \sim \pi_{\theta_{old}} (\cdot \mid q)$
      \State Evaluate each $o_i$ using $r_{\varphi}$ to obtain rewards $\{r_i\}_{i=1}^{G}$
      \State Use group-based estimation to compute $\hat{A}_{i,t}$ for token $t$ in $o_i$
      \For{GRPO iteration = 1, \dots, $\mu$}
        \State Update policy $\pi_{\theta}$ by maximizing the GRPO objective (Equation \ref{eq:GC-GRPO})
      \EndFor
    \EndFor \State Periodically retrain $r_\varphi$ with a replay buffer. 
    \EndFor 
  \end{algorithmic}
  \textbf{Output} $\pi_\theta$
  \label{alg:iter-grpo}
\end{algorithm}

\subsubsection{Outcome Supervision RL with GRPO} In the outcome supervision framework, for every question $q$, a set of $G$ output samples $\{o_1, o_2, \cdots, o_G\}$ is drawn from the policy model snapshot $\pi_{\theta_{old}}$. The reward model scores each output, providing a reward vector $\mathbf{r}=\{r_1, r_2, \cdots, r_G\}$ aligned with the outputs. These reward values are standardized within the group by removing the mean and dividing by the standard deviation. For each output $o_i$, the normalized reward is then used as a constant advantage $\hat{A}_{i, t}$ across every token, specifically $\hat{A}_{i, t} = \widetilde{r}_i = \frac{r_i- {\rm mean}(\mathbf{r})}{{\rm std}(\mathbf{r})}$. The policy parameters are then adjusted by maximizing the objective presented in equation (\ref{eq:GRPO-obj}).

\subsubsection{Process Supervision RL with GRPO} Relying exclusively on terminal rewards may not adequately guide learning, especially for tasks involving intricate reasoning as in mathematics. Building upon approaches in \cite{wang2023math}, we consider process supervision, in which the reward is given after each intermediate reasoning step. Concretely, for a question $q$ and $G$ sampled completions $\{o_1, o_2, \cdots, o_G\}$, a process reward model assigns rewards to each reasoning step, constructing a set $\mathbf{R} = \{ \{r_1^{index(1)},\cdots,r_1^{index(K_1)}\}, \cdots,  \{r_G^{index(1)},\cdots,r_G^{index(K_G)}\} \}$, where $index(j)$ denotes the final token index of the $j$-th step and $K_i$ is the step count for $o_i$. These step-wise rewards are similarly normalized: $\widetilde{r}_i^{index(j)} = \frac{r_i^{index(j)} - {\rm mean(\mathbf{R})}}{{\rm std(\mathbf{R})}}$.
Process supervision then defines the advantage for each token $t$ in $o_i$ as the sum of all normalized future step rewards, i.e., $\hat{A}_{i, t} = \sum_{index(j) \ge t} \widetilde{r}_i^{index(j)}$. Optimization proceeds by maximizing the same objective in equation (\ref{eq:GRPO-obj}).

\subsubsection{Iterative RL with GRPO} During the course of reinforcement learning, earlier versions of the reward model may not provide adequate supervision for the evolving policy model. To address this limitation, we investigate the use of iterative RL with GRPO. As detailed in Algorithm \ref{alg:iter-grpo}, our approach involves periodically regenerating the reward model’s training dataset using samples drawn from the current policy model. The previous reward model is continually updated via a replay mechanism, which incorporates 10\% of data from earlier training iterations. Subsequently, the updated policy model becomes the new reference model, and training proceeds with this policy model under the guidance of the newly trained reward model.

\subsection{Training and Evaluating DeepSeekMath-RL}

Our reinforcement learning experiments leverage the DeepSeekMath-Instruct 7B model. The RL training corpus comprises approximately 144K chain-of-thought questions sourced from GSM8K and MATH datasets in the SFT collection. We deliberately exclude other SFT instances to assess how RL affects benchmarks when their data are not present in the RL training set. The reward model’s training data are prepared in line with the procedure described by \citep{wang2023math}. We initialize training of the reward model on DeepSeekMath-Base 7B with a learning rate of 2e-5. For GRPO, the learning rate assigned to the policy model is 1e-6, and the KL regularization term is set to 0.04. For each problem, we generate $64$ samples. The sequence length is capped at 1024, and training is performed in batches of 1024. The policy model is updated only once after each exploration phase. For evaluation, we assess DeepSeekMath-RL 7B using the same set of benchmarks as DeepSeekMath-Instruct 7B. In this context, chain-of-thought versions of GSM8K and MATH represent in-domain benchmarks for DeepSeekMath-RL 7B, whereas all other evaluation tasks are treated as out-of-domain.

Table \ref{tab:sft_rl_math} presents a comparative analysis of the performance of both open- and closed-source models that utilize chain-of-thought and tool-integrated reasoning strategies, evaluated across English and Chinese benchmarks. The key observations are as follows: 1) When employing chain-of-thought reasoning, DeepSeekMath-RL 7B achieves accuracies of 88.2\% on GSM8K and 51.7\% on MATH, outperforming all open-source competitors within the 7B to 70B parameter range and exceeding the performance of most closed-source systems. 2) Notably, DeepSeekMath-RL 7B's training is limited solely to chain-of-thought-format instruction data from GSM8K and MATH, using DeepSeekMath-Instruct 7B as its initialization. Even with such a focused training regime, DeepSeekMath-RL 7B demonstrates superior results across every evaluation metric compared to DeepSeekMath-Instruct 7B, thereby highlighting the impact and efficacy of reinforcement learning.

\section{Discussion}
This section provides an overview of our insights from both pre-training and reinforcement learning experiments.

\subsection{Lessons Learnt in Pre-Training}

We begin by discussing our experience during the pre-training phase. Unless noted otherwise, our experiments utilize the training protocol described in Section \ref{sec:quality-policy}. It should be emphasized that references to the DeepSeekMath Corpus here pertain to an 89B-token collection assembled during the second round of data gathering.

\subsubsection{Code Training Benefits Mathematical Reasoning}

There is a widely circulated but not yet substantiated hypothesis suggesting that training on code can boost reasoning ability. Here, we offer partial empirical evidence supporting this claim in the context of mathematical problem-solving: code pre-training leads to significant improvements in models’ mathematical reasoning capabilities, with or without access to external tools.

To assess the influence of code training on mathematical reasoning, we conducted investigations under two experimental paradigms: two-stage and one-stage training.

\textbf{Two-Stage Training}

\begin{itemize}[topsep=0pt]
    \item \textbf{Code Training for 400B Tokens $\rightarrow$ Math Training for 150B Tokens}: DeepSeek-LLM 1.3B is subjected to pretraining on 400B code tokens, which is then followed by further training on 150B math tokens.
    \item \textbf{General Training for 400B Tokens $\rightarrow$ Math Training for 150B Tokens}: For comparative purposes, another model is trained first on 400B general tokens (sampled from DeepSeek-AI’s comprehensive general corpus) instead of code tokens, before undergoing math training. This setup serves to explore whether pretraining on code offers particular advantages over general-domain text in terms of enhancing mathematical reasoning capabilities.
\end{itemize}

\textbf{One-Stage Training}

\begin{itemize}[topsep=0pt]
    \item \textbf{Math Training for 150B Tokens}: The model is exclusively trained on 150B math tokens without any prior code pretraining.
    \item \textbf{Training on a mixture of 400B Code Tokens and 150B Math Tokens}: Given that conducting math training after code pretraining tends to deteriorate coding ability, we test a one-stage mixed-token approach. Here, code and math tokens are presented together during training, aiming to support mathematical reasoning while also preventing catastrophic forgetting of code-related capabilities.
\end{itemize}

\paragraph{Results}
Tables \ref{tab:code-math} and \ref{tab:code-math-general-eval} present the comparative downstream results for these training paradigms.

Incorporating code pretraining enhances program-aided mathematical reasoning, regardless of whether the approach involves sequential or joint training stages. According to Table \ref{tab:code-math}, code pretraining alone markedly elevates the model’s performance on GSM8K and MATH benchmarks with Python-based solutions. Subsequent math-specific training further boosts these outcomes. Notably, simultaneous exposure to both code and math tokens in a single training phase not only reduces catastrophic forgetting associated with staged training, but also fosters proficiency in coding (Table \ref{tab:code-math-general-eval}) as well as in program-supported mathematical reasoning (Table \ref{tab:code-math}).

Furthermore, code pretraining proves beneficial for mathematical reasoning even without reliance on external tools. Within the two-stage scheme, the initial code-focused training already imparts noticeable improvements, and it facilitates more efficient later math-specific learning, culminating in optimal performance. However, concurrently training on mixed code and math tokens in a single stage negatively affects non-tool-based mathematical reasoning. One possible explanation is that DeepSeek-LLM 1.3B, owing to its relatively modest scale, may not possess sufficient capacity to optimally integrate code and mathematics knowledge concurrently.

\subsubsection{ArXiv Papers Appear Inefficient in Enhancing Mathematical Reasoning}

ArXiv papers are frequently utilized as a part of mathematical pre-training datasets \citep{minerva,gpt-f,llemma,mathpile}. Nevertheless, there has been a lack of comprehensive studies that rigorously assess their contribution to mathematical reasoning abilities. Somewhat unexpectedly, our experimental findings suggest that arXiv papers contribute little, if at all, to advances in mathematical reasoning. We evaluated models at various scales, such as DeepSeek-LLM 1.3B and DeepSeek-Coder-Base-v1.5 7B \citep{deepseek-coder}, with the arXiv data being processed through distinct data curation procedures:

\begin{itemize}[topsep=0pt]
    \item \textbf{MathPile} \citep{mathpile}: a dataset containing 8.9B tokens, assembled through cleaning and heuristic filtering, with over 85\% of the content derived from scientific arXiv articles;
    \item \textbf{ArXiv-RedPajama} \citep{redpajama}: this resource comprises all arXiv LaTeX documents, stripped of preambles, commentary, macros, and references, comprising 28.0B tokens in total.
\end{itemize}

We conducted isolated pre-training runs: DeepSeek-LLM 1.3B was trained for 150B tokens, and DeepSeek-Coder-Base-v1.5 7B was trained for 40B tokens, each exclusively on an arXiv-based corpus. The results indicate that, for both models, reliance on arXiv-only pre-training fails to deliver discernible enhancements and can even undermine performance on a variety of mathematical tasks with differing levels of complexity. Our benchmarks span quantitative reasoning (e.g., GSM8K and MATH, see Table \ref{tab:arxiv-cot}), multiple-choice STEM evaluations such as MMLU-STEM (Table \ref{tab:arxiv-cot}), as well as formal mathematics, e.g., miniF2F (Table \ref{tab:arxiv_atp}).

Nonetheless, these findings have caveats that merit consideration. Our current analysis does not cover:

\begin{itemize}[topsep=0pt]
    \item How arXiv-derived tokens might affect other mathematics-specific tasks absent from this evaluation, for instance, the informalization of theorems, where formal statements or proofs are rewritten in informal language;
    \item The role of arXiv data when integrated alongside other modalities or datasets;
    \item The possibility that arXiv content might be more advantageous at greater model scales.
\end{itemize}

As such, these questions remain open for subsequent investigation.

\subsection{Reflections on Reinforcement Learning}
\subsubsection{Towards a Unified Training Framework} This portion proposes an overarching framework for analyzing a range of training techniques, including SFT, RFT, DPO, PPO, and GRPO, followed by empirical studies examining elements of the proposed framework. In general terms, for these methods, the gradient of the loss function with respect to model parameters $\theta$ can be represented as:

\begin{equation}
    \nabla_{\theta}\mathcal{J}_{\textcolor{red}{\mathcal{A}}}(\theta) = \mathbb{E}[\underbrace{(q,o) \sim \textcolor{red}{\mathcal{D}}}_{Data \ Source}]\left( \frac{1}{|o|} \sum_{t=1}^{|o|}  \underbrace{GC_{{\mathcal{A}}}(q, o, t, \textcolor{red}{\pi_{{rf}}})}_{Gradient \ Coefficient}  \nabla_{\theta}\log \pi_{\theta}(o_t | q, o_{<t})\right).
\label{eq:objective}
\end{equation}

This framework comprises three essential elements: 1) \textit{Data Source $\mathcal{D}$}, which specifies the set of samples used for learning; 2) \textit{Reward Function $\pi_{{rf}}$}, providing the evaluative feedback that guides training; 3) \textit{Algorithm $\mathcal{A}$}, responsible for translating data instances and reward signals into the gradient coefficient $GC$, which quantifies how strongly each sample influences the parameter updates—either by penalization or reinforcement. We compare several notable training techniques within this overarching scheme:

\begin{itemize}
    \item  \textbf{Supervised Fine-tuning (SFT)}: SFT adapts a pretrained model by training on curated human-labeled SFT datasets.
    \item \textbf{Rejection Sampling Fine-tuning (RFT)}: RFT continues to refine an SFT-initialized model using outputs produced by the SFT model itself, keeping only those answers to SFT prompts that meet predefined correctness criteria.
    \item \textbf{Direct Preference Optimization (DPO)}: DPO takes the SFT model and fine-tunes it further on augmented datasets produced from SFT outputs, employing a pairwise DPO loss to encode preference judgments.
    \item \textbf{Online Rejection Sampling Fine-tuning (Online RFT)}: Unlike standard RFT, Online RFT starts with the SFT model but draws augmented outputs directly from the evolving, real-time policy model, then uses these outputs for additional fine-tuning.
    \item \textbf{PPO/GRPO}: PPO/GRPO use the SFT model as the policy initialization and optimize it further through reinforcement, leveraging outputs sampled from the current, dynamically updated policy model.
\end{itemize}

We provide an overview of these methods' components in Table \ref{tab:data_policy}. For a comprehensive explanation of the derivation steps, please consult Appendix \ref{app:analysis-rl}.

\paragraph{Observation about Data Source} The data sources are categorized into two types: online sampling and offline sampling. Online sampling refers to acquiring training data through exploration performed by the current, real-time policy model during training. In contrast, offline sampling utilizes training data generated from the output of the initial SFT model. Methods such as RFT and DPO employ the offline sampling paradigm, whereas Online RFT and GRPO rely on the online sampling approach.

As depicted in Figure \ref{fig:rl-analysis}, our results indicate that Online RFT achieves markedly superior performance over RFT across both benchmarks. While both methods perform comparably at the outset of training, Online RFT establishes a clear and growing performance lead during later training stages, underscoring the benefits of an online sampling strategy. This observation aligns with expectations: at the beginning of training, the behavior of the actor model closely mirrors that of the SFT model, resulting in sampled data with only slight differences. As training progresses, the actor's output diverges more from the SFT model, and sampling directly from the actor provides increasingly informative training data, thus amplifying the advantages of online sampling.

\paragraph{Observation about Gradient Coefficient} In these algorithms, the reward signal is transformed into a gradient coefficient, which in turn drives model parameter updates. We categorize the reward function into `Rule' and `Model' in our investigation. The `Rule' category assesses the appropriateness of responses by evaluating answer correctness, whereas the `Model' approach involves training a reward model to assign a score to each response, with the reward model's own training data grounded in rule-based assessments. Equations \ref{eq:GC-RFT} and \ref{eq:GC-GRPO} elucidate an essential distinction between GRPO and Online RFT: GRPO modifies its gradient coefficient in response to the reward model's outputs, allowing it to apply more granular reinforcement or penalty depending on the response's quality. Conversely, Online RFT does not differentiate in this way—it treats all correct answers uniformly, failing to penalize incorrect answers and providing the same degree of positive reinforcement to all correct responses.

As illustrated in Figure \ref{fig:rl-analysis}, GRPO achieves better results than online RFT, thereby underscoring the effectiveness of modifying both positive and negative gradient coefficients. Moreover, GRPO+PS consistently outperforms GRPO+OS, highlighting the advantage of employing gradient coefficients that are step-aware and finely tuned. In addition, we examine the application of iterative RL in our experiments, conducting two consecutive rounds. According to the results presented in Figure \ref{fig:iter-rl}, the introduction of iterative RL yields marked gains, particularly during the initial iteration.

\subsubsection{Why RL Works?} In this study, reinforcement learning is conducted using a selected subset of instruction tuning data, resulting in a notable boost over the baseline instruction tuning model. To further clarify the mechanisms underlying reinforcement learning's effectiveness, we assess both Pass@K and Maj@K accuracy for the Instruct and RL-enhanced models across two standard benchmarks. The outcomes, visualized in Figure \ref{fig:maj-pass}, show that reinforcement learning advances Maj@K but has negligible impact on Pass@K. This pattern implies that the overall improvement conferred by reinforcement learning is largely due to increased robustness of the model’s output distribution; that is, \textbf{it appears the primary contribution is lifting the rank of correct answers within the TopK set, rather than improving intrinsic problem-solving abilities.} Similarly, \citep{wang2023making} observe a \textbf{misalignment problem} in reasoning benchmarks when using SFT models, and find that various preference alignment methodologies \citep{yuan2023rrhf,song2023preference,wang2023making} can lead to substantial gains in reasoning capacity for SFT systems.

\subsubsection{How to Achieve More Effective RL?}
\label{sec:effec_rl}
We have shown that reinforcement learning is particularly effective for mathematical reasoning problems, and we additionally introduce a unified perspective for interpreting various representative training techniques. Within this framework, different approaches may be viewed as forms or special cases of RL-based methods. According to Equation \ref{eq:objective}, three central elements are identified: Data Source, Algorithm, and Reward Function. We outline potential directions for advancement for each of these components.

\paragraph{Data Source} The data source constitutes the foundational input for all training paradigms. For reinforcement learning scenarios, we specifically define the data source as the set of unlabeled questions alongside policy model-generated responses. In this work, we restrict our exploration to questions drawn from the instruction tuning phase, using basic nucleus sampling to produce outputs. We hypothesize this is a key factor behind our RL pipeline’s improvements being mostly limited to Maj@K. In future work, we plan to expand our RL procedure to encompass out-of-distribution prompts, leveraging \textbf{more sophisticated sampling (decoding) strategies}, such as those based on tree-search algorithms \citep{yao2023tree}. Additionally, incorporating \textbf{advanced inference methodologies} \citep{xia-etal-2023-speculative,leviathan2023fast,kwon2023efficient,xia2024unlocking}—which are crucial for efficient policy model exploration—remains a vital area for enhancing performance.

\paragraph{Algorithms} Algorithms leverage both the input data and the associated reward signal to adjust model parameters by modifying the gradient coefficient. As outlined in Equation \ref{eq:objective}, contemporary approaches largely place complete \textbf{TRUST} in the reward function’s feedback, modulating the conditional probability of particular tokens accordingly. Nevertheless, the reliability of the reward signal cannot be guaranteed, particularly for highly intricate tasks. For instance, the PRM800K datasets \citep{lightman2023let}, despite being meticulously labeled by skilled annotators, still exhibit about 20\% incorrect labels\footnote{\url{https://github.com/openai/prm800k/issues/12\#issuecomment-1728491852}}. Consequently, we aim to investigate reinforcement learning algorithms that demonstrate robustness in the presence of noisy or unreliable reward signals. We anticipate that alignment strategies exemplified by \textbf{WEAK-TO-STRONG} methods \citep{burns2023weak} will substantially reshape the foundations of current learning algorithms.

\paragraph{Reward Function} The reward function serves as the primary source of learning signals during training. In the context of RL, this role is typically fulfilled by a neural reward model. We identify three pivotal challenges for advancing reward models: 1) \textbf{Improving the generalization capacity of reward models.} Reward models must be capable of reliably evaluating out-of-distribution queries and sophisticated outputs; failing this, reinforcement learning may simply consolidate the distribution of large language models (LLMs) without yielding genuine progress in their abilities; 2) \textbf{Quantifying and leveraging reward model uncertainty.} Capturing the uncertainty associated with reward models may provide a critical link between underperforming reward models and weak-to-strong alignment approaches; 3) \textbf{Developing high-fidelity, process-level reward models efficiently} to deliver detailed training feedback during reasoning tasks \citep{lightman2023let,wang2023math}.

\section{Conclusion, Limitation, and Future Work} 
We introduce DeepSeekMath, which establishes a new benchmark by surpassing all available open-source models on the competitive MATH dataset and nearly matches the results of leading proprietary models. DeepSeekMath~is initialized from DeepSeek-Coder-v1.5 7B and is trained continuously for 500B tokens, with 120B of those being math-specific tokens extracted from Common Crawl making up a substantial proportion of the data. Our comprehensive ablation analyses highlight that web data can serve as a rich resource for acquiring high-quality mathematical datasets, whereas arXiv content may not provide the anticipated advantages. We propose Group Relative Policy Optimization (GRPO), an adaptation of Proximal Policy Optimization (PPO), which demonstrably enhances mathematical reasoning skills while consuming less memory. Experimental evidence confirms that GRPO yields meaningful gains even when DeepSeekMath-Instruct 7B already performs strongly on evaluation suites. Additionally, we outline a cohesive framework for interpreting various techniques and discuss several promising avenues to achieve more efficient reinforcement learning.

Despite DeepSeekMath's~strong results on quantitative reasoning assessments, it remains less proficient than closed models in areas such as geometry and theorem-proving. For example, in preliminary tests, the system struggled with problems concerning triangles and ellipses, pointing to a potential bias in the chosen pre-training and fine-tuning data. Furthermore, limitations in model size contribute to DeepSeekMath's~inferior few-shot learning capability compared to GPT-4, which shows marked improvements when few-shot prompts are used, whereas DeepSeekMath~displays comparable performance in both zero-shot and few-shot contexts. Looking forward, our plan is to further optimize our data engineering pipeline in order to curate even higher-quality pre-training corpora. We also intend to investigate additional strategies (Section \ref{sec:effec_rl}) to advance reinforcement learning methods for large language models.